\section{Method}
\label{sec:method}

We begin with preliminaries and notation (Sec.~\ref{sec:method:prelim}), define the shell entropy concentration $\rho_f(\tau)$ and report its value on the checkpoints we study (Sec.~\ref{sec:method:rho}), state the Post-hoc Entropy Minimization loss (Sec.~\ref{sec:method:loss}), and give the intuition for why a converged checkpoint is the right regime for it (Sec.~\ref{sec:method:intuition}). The formal argument for why the loss is well-conditioned in this regime is deferred to Section~\ref{sec:theory}.

\subsection{Preliminaries and notation}
\label{sec:method:prelim}

Let $\mathcal{U}$ denote the set of unlabeled training volumes and $\mathcal{L}$ the set of labeled training volumes, both drawn from the same distribution. Let $f_{\theta^*}: \mathbb{R}^{H\times W\times D} \to \mathbb{R}^{C \times H \times W \times D}$ denote a segmentation network converged under an SSL objective on $\mathcal{L} \cup \mathcal{U}$; in our experiments $f_{\theta^*}$ is a V-Net~\cite{milletari2016vnet} trained with BCP~\cite{bai2023bcp}, but the analysis that follows does not depend on this choice. For a volume $\mathbf{x} \in \mathcal{U}$ and a voxel index $i$, write
\begin{equation}
p_i = \mathrm{softmax}\big(f_{\theta^*}(\mathbf{x})_i\big) \in \Delta^{C-1},
\qquad
H_i = -\sum_{c=1}^{C} p_{i,c}\,\log p_{i,c},
\qquad
c_i = \max_{c} p_{i,c},
\end{equation}
for the per-voxel softmax distribution, per-voxel Shannon entropy, and per-voxel confidence respectively. We will refer to voxels with $c_i\!>\!\tau$ for some threshold $\tau$ close to $1$ as \emph{saturated} and voxels with $c_i\!\le\!\tau$ as \emph{uncertain}. The total prediction entropy of the network on a volume is $H(\mathbf{x}) = \sum_i H_i(\mathbf{x})$, and its expectation over $\mathcal{U}$ is $\bar H = \mathbb{E}_{\mathbf{x}\sim\mathcal{U}}[H(\mathbf{x})]$.

\subsection{Shell entropy concentration $\rho_f(\tau)$}
\label{sec:method:rho}

The observation that motivates PEM is not about the magnitude of residual entropy on unlabeled data, but about its spatial distribution. Two networks can have similar total $\bar H$ while differing sharply in \emph{where} that entropy lives: one network may carry its residual uncertainty on a thin boundary shell where the underlying evidence is genuinely ambiguous, while another may have diffuse uncertainty scattered across interior regions the network does not yet understand. These two networks are in qualitatively different regimes, and an entropy-minimization loss behaves qualitatively differently on them. We capture this distinction with a single scalar.

\begin{definition}[Shell entropy concentration]
\label{def:rho}
Given a converged segmentation network $f_{\theta^*}$, an unlabeled volume distribution $\mathcal{U}$, and a confidence threshold $\tau \in (0, 1]$, define
\begin{equation}
\rho_f(\tau) \;=\; \frac{\mathbb{E}_{\mathbf{x}\sim\mathcal{U}}\!\left[\sum_i H_i(\mathbf{x})\,\mathbb{1}[c_i(\mathbf{x}) \le \tau]\right]}{\mathbb{E}_{\mathbf{x}\sim\mathcal{U}}\!\left[\sum_i H_i(\mathbf{x})\right]} \;\in\; [0, 1].
\label{eq:rho}
\end{equation}
\end{definition}

$\rho_f(\tau)$ is the fraction of the network's total prediction entropy that is carried by voxels of confidence at most $\tau$. We refer to the set of such voxels as the \emph{uncertain shell} of the network at level $\tau$; for a well-converged segmentation network, this shell is spatially localized near the network's predicted decision surface, and for this reason it often consists of voxels lying on or near the physical organ boundary (we verify this empirically below, and explain why formally in Sec.~\ref{sec:theory}).

\begin{proposition}[Basic properties of $\rho_f(\tau)$]
The quantity $\rho_f(\tau)$ is: \emph{(i)} label-free, since it depends only on forward passes of $f_{\theta^*}$ over $\mathcal{U}$; \emph{(ii)} monotone non-decreasing in $\tau$, with $\rho_f(0)\!=\!0$ and $\rho_f(1)\!=\!1$; \emph{(iii)} scale-invariant with respect to the total residual entropy $\bar H$, so that two networks with very different $\bar H$ can be directly compared; and \emph{(iv)} scalar, making it a convenient summary statistic to report alongside any released SSL checkpoint.
\end{proposition}

Property (iii) is the one that makes $\rho_f$ useful for the argument we want to make. The absolute magnitude of residual entropy in a converged segmentation network depends strongly on the network architecture, the training loss, the calibration temperature of the final layer, and the amount of augmentation --- none of which is directly related to whether the network is ready for a refinement step. $\rho_f(\tau)$ factors these out: a highly confident network and a mildly confident network can both have $\rho_f(\tau)$ close to $1$, as long as whatever residual entropy they carry is concentrated in the uncertain shell.

\paragraph{A PEM-suitability diagnostic.} Throughout this paper we report $\rho_f$ at $\tau\!=\!0.99$, and we use the empirically determined threshold $\rho^\star\!=\!0.85$ as a suitability check: we expect PEM to be well-conditioned on any checkpoint for which $\rho_f(0.99) \ge \rho^\star$, and ill-conditioned otherwise. The choice of $\tau\!=\!0.99$ is not essential to the definition, and is made for two reasons. First, it is conservative: it excludes from the uncertain shell only the most saturated voxels, so the shell contains essentially all voxels where the network is not already making a very confident prediction. Second, it is the level at which the distinction between SSL-converged and supervised-only checkpoints is sharpest in our measurements (Table~\ref{tab:entropy}); weaker thresholds dilute the separation, and stronger thresholds leave too few voxels in the shell for the statistics to be stable.

\paragraph{Measured shell concentration on the checkpoints we study.} We measure $\rho_f(0.99)$ on the publicly released BCP checkpoints for Pancreas-CT 20\%, LA 5\%, and LA 10\%, and on a supervised-only V-Net trained on the $12$ Pancreas-CT labeled volumes from the BCP split. The three SSL checkpoints concentrate $\rho_f(0.99)\!\in\!\{0.894, 0.948, 0.945\}$ respectively. The supervised-only baseline concentrates only $\rho_f(0.99)\!=\!0.552$: nearly half of its residual entropy is spread across interior regions rather than the boundary shell. To quantify how far off the predicted boundary the uncertain voxels lie, we compute the mean $\ell_2$ distance from each uncertain voxel to the nearest predicted decision-surface voxel; this distance is $5$--$11$ voxels on the SSL checkpoints and $14.5$ voxels on the supervised baseline. The gap between the two regimes is a $34$ to $40$ percentage-point difference in $\rho_f$ plus a roughly $3\times$ difference in uncertain-voxel-to-boundary distance, and it is observed consistently across two datasets and three label fractions with no tuning. These numbers are reported in Table~\ref{tab:entropy}.

\subsection{The PEM loss}
\label{sec:method:loss}

Post-hoc Entropy Minimization takes a converged checkpoint $f_{\theta^*}$ and continues optimization for a small fixed number of additional epochs against the following loss:
\begin{equation}
\mathcal{L}_{\mathrm{PEM}}(\theta; \mathcal{U}) \;=\; \frac{1}{|\mathcal{U}|} \sum_{\mathbf{x} \in \mathcal{U}} \frac{\sum_i m_i(\mathbf{x})\, H_i(\mathbf{x}; \theta)}{\sum_i m_i(\mathbf{x}) + \epsilon},
\qquad
m_i(\mathbf{x}) \;=\; \mathbb{1}\!\big[c_i(\mathbf{x}; \theta^*) \le \tau\big],
\label{eq:pem}
\end{equation}
where $\epsilon > 0$ is a small stability constant, $H_i$ is computed with respect to the current fine-tuning parameters $\theta$, and the mask $m_i$ is computed from the \emph{initial} converged checkpoint $\theta^*$ and held fixed throughout the fine-tune. This last choice is deliberate: recomputing the mask at every step would make the loss nonstationary in a way that couples the gradient at step $t$ to its own previous updates, and in practice causes the mask to shrink as the network sharpens its predictions, amplifying noise. Freezing the mask at initialization keeps the loss stationary and the optimization convex in the neighborhood of $\theta^*$. Setting $\tau\!=\!1$ recovers full-volume entropy minimization (all voxels weighted equally, no explicit shell restriction); setting $\tau\!<\!1$ restricts the loss to the uncertain shell defined by the initial checkpoint. We report results for both cases in the ablations (Sec.~\ref{sec:ablations}).

The loss in Eq.~\ref{eq:pem} is, formally, identical to the entropy term of Grandvalet and Bengio~\cite{grandvalet2004entropy}, masked by a confidence gate. What is non-classical about PEM is not the object we are minimizing but the state of the network at which we begin minimizing it. Section~\ref{sec:theory} explains why this difference is consequential.

\subsection{Why a converged checkpoint is the right regime}
\label{sec:method:intuition}

It is tempting to view an optimization loss as a static object whose properties are independent of the state of the network it is applied to. For many losses this view is adequate; for entropy minimization it is misleading. The central intuition behind PEM is that the gradient of the entropy loss is only as useful as the state of the posterior it acts on, and that a converged SSL checkpoint and a supervised-only baseline are in qualitatively different posterior states.

\paragraph{Diffuse entropy is noise; concentrated entropy is signal.}
Consider two networks with the same total residual entropy $\bar H$ on an unlabeled volume. The first is a converged SSL network whose entropy is carried by $5\%$ of the voxels, all of them within a few voxels of the predicted decision surface. The second is an undertrained or supervised-only network whose entropy is scattered across $50\%$ of the voxels, half of them in interior regions where the network has not learned tissue texture and the other half on the boundary. The entropy gradient $-\nabla_\theta H(p_i)$ for both networks has the same magnitude, but its \emph{directional coherence} across voxels differs sharply. In the first case, the gradient at every active voxel points in a direction that refines the nearby boundary --- there is work to be done and the work is local. In the second case, the gradient at interior voxels commits the network to local minima it does not yet understand, and the gradient at boundary voxels competes with the interior gradients for the limited step-size budget of a gradient-descent update. The optimization step that minimizes the sum of all these gradients is, in expectation, a smaller and noisier version of the step that the first network would have taken if its boundary voxels had been acted on in isolation.

This is the root reason why training-time entropy minimization is a delicate object: it needs careful balancing against a supervised loss and typically some form of class-balance or curriculum regularization to prevent collapse, because the network distributions it acts on are not yet concentrated. Post-hoc entropy minimization on a converged SSL checkpoint does not need this stabilizing machinery because the concentration has already happened --- the long SSL training run has already sharpened the interior of each organ and left only the boundary as a source of residual uncertainty. The post-hoc regime inherits the convergence of the SSL training and pays for it only once, rather than reinventing a balancing problem at every step.

\paragraph{Why concentration emerges at convergence.}
Why do converged SSL networks concentrate residual entropy at organ boundaries in the first place? The answer, at a high level, is that the labeled fraction of the training set already determines the interior of each organ with high confidence. For a voxel lying well inside the pancreas in a labeled volume, the supervised loss has seen examples of the class label and the local tissue texture; the network learns to recognize the texture and commits to the interior class early in training. This commitment propagates to unlabeled volumes via the consistency or copy-paste mechanism of the SSL method. Voxels near the physical organ surface, however, are harder: their feature distribution is inherently bimodal because partial-volume averaging, scanner blur, and biological variability all cause the same physical voxel to contain signal from two tissues. No amount of supervised evidence can fully resolve this bimodality on a finite labeled set, and the SSL training process therefore leaves exactly these voxels as the residual source of uncertainty. In our experiments, the BCP checkpoints carry $89$--$95\%$ of their residual entropy within $5$--$11$ voxels of the predicted organ surface, and only $5$--$9\%$ of the volume is ever above this uncertain-shell threshold. The supervised-only baseline, which does not benefit from the SSL consistency mechanism propagating the interior commitment to unlabeled data, carries only $55\%$ of its entropy in the corresponding shell and places its uncertain voxels on average $14.5$ voxels from the boundary --- far enough that the entropy-minimization gradient, if applied to them, would be acting on voxels that are not near any decision surface at all.

\paragraph{The $p(1{-}p)$ factor as a preview.}
A full formal account of why this matters is given in Section~\ref{sec:theory}, where we derive the binary-entropy gradient $\partial H/\partial z\!=\!-z\,p(1{-}p)$ and show that it acts as an implicit boundary-shell gate: the update vanishes on saturated voxels ($p$ close to $0$ or $1$), peaks on boundary voxels ($p\!\approx\!0.5$), and therefore does useful work only where there is work to be done. Combined with the convergence-induced concentration of $\rho_f$, this implicit gating explains both why PEM does not need an explicit class-balance regularizer to prevent collapse and why it shrinks in magnitude as the base checkpoint becomes more saturated (the second failure mode we observe on LA 10\%; Sec.~\ref{sec:discussion}). For the present section it suffices to say: the reason post-hoc entropy minimization works where the training-time version needs extensive stabilization is that the convergence of the base model has done the stabilization for us.

\subsection{Implementation and fine-tuning protocol}
\label{sec:method:impl}

We fine-tune the publicly released BCP checkpoints for Pancreas-CT 20\% (V-Net with instance normalization, $12$ labeled $/$ $50$ unlabeled $/$ $18$ test) and LA 2018 (V-Net with batch normalization, $4$ or $8$ labeled $/$ rest unlabeled $/$ $20$ test) using BCP's canonical preprocessing and data splits. Our reproduction of the BCP baseline matches the reported Dice to within $0.02$ on all three configurations. Training patches are $96^3$ on Pancreas-CT and $112\!\times\!112\!\times\!80$ on LA, center-cropped from the preprocessed H5 volumes with batch size $B\!=\!2$ and a seed-dependent shuffled loader; we use no flipping or random cropping to match the BCP evaluation protocol. The optimizer is Adam with no weight decay, no learning-rate schedule, and $\ell_2$ gradient clipping at norm $1.0$. The learning rate is chosen once per dataset via $5$-fold cross-validation on the labeled training volumes --- $5\!\times\!10^{-5}$ for Pancreas-CT and $1\!\times\!10^{-5}$ for LA, the latter being gentler because the LA base checkpoint is more saturated and admits a smaller step size before the optimization begins to relocate rather than refine the decision boundary. For LA's BatchNorm backbone we freeze all running statistics and only update the affine scale and shift parameters; for Pancreas-CT's InstanceNorm backbone no such freezing is needed. The number of PEM epochs is fixed in advance at $E\!=\!2$ for Pancreas-CT and $E\!=\!5$ for LA, for a total of $50$, $190$, and $180$ parameter updates per run on the three configurations respectively; no quantity is read from the test set during fine-tuning, and no early stopping is used beyond a collapse-protection safeguard that halts the run if the training loss increases by more than $0.5\%$ from its initial value (triggered only in diagnosed-ill-conditioned regimes; see Sec.~\ref{sec:discussion}). Inference follows BCP's original sliding-window protocol with patch size equal to the training patch and stride $(16, 16, 4)$ on Pancreas-CT and $(18, 18, 4)$ on LA, no test-time augmentation. The three-seed PEM rows of Table~\ref{tab:main} are averaged over seeds $2020$, $42$, $123$; all other post-hoc baselines and ablations use seed $2020$. Per-configuration wall-clock on a single NVIDIA RTX $4070$ Ti SUPER is $5.1$ minutes for Pancreas-CT 20\%, $2.7$ minutes for LA 5\%, and $2.6$ minutes for LA 10\%; the full experimental grid --- three-seed PEM, three controlled post-hoc baselines, and the learning-rate $\times$ threshold $\times$ epoch-budget ablation on LA 5\% --- completes in under four GPU-hours.

\paragraph{Determinism and reproducibility.} All runs use \texttt{torch.use\_deterministic\_algorithms(True)}, \texttt{cudnn.deterministic=True}, \texttt{cudnn.benchmark=False}, and a seeded random number generator for every stochastic component (weight initialization, data loader shuffling, augmentations). Fixed seeds produce bit-exact reproductions of our reported numbers on the hardware we used. Code, released checkpoints, split files, per-run logs, and the exact frozen configuration of every experiment in the paper will be released at \url{<anonymized>} upon acceptance.
